\section{Introduzione}\label{sec:Introduzione}

Le password sono ancora il principale meccanismo di autenticazione dei servizi
informatici, dai social network all'e-commerce: a differenza di token hardware o
sistemi biometrici, non richiedono dispositivi dedicati e sono comode per l'utente.
Questa stessa semplicità le rende però il punto debole del sistema: il \emph{Data Breach Investigations Report} di
Verizon~\cite{verizon2024dbir} documenta che l'uso di credenziali rubate o
deboli è il punto d'ingresso più frequente nelle violazioni informatiche.

Il problema di fondo è che le password non vengono scelte in modo uniforme.
Bonneau~\cite{bonneau2012}, analizzando un dataset anonimizzato di 70 milioni di
password Yahoo, ha mostrato che la distribuzione reale delle scelte umane è fortemente
sbilanciata verso pochi schemi ricorrenti, e lo studio su larga scala di
Flor\^{e}ncio e Herley~\cite{florencio2007}, condotto su circa 500.000 utenti per
tre mesi, ha documentato i comportamenti che causano questo squilibrio, a partire
dal riuso delle password e dalla preferenza per combinazioni facili da memorizzare.
È proprio su queste abitudini che si fondano gli attacchi a dizionario.

Modellando la struttura delle password con grammatiche probabilistiche (PCFG)~\cite{weir2009}
si possono generare candidati in ordine decrescente di probabilità, con tassi di
successo molto più alti rispetto al brute-force o alle wordlist statiche.
Anche le liste più estese (come quelle derivate dal breach RockYou) restano però generiche e poco efficaci quando l'utente introduce
variazioni personali.

Il salto di qualità successivo sono i dizionari mirati sulla singola vittima:
il framework TarGuess~\cite{wang2016targuess} ne automatizza la generazione estraendo
dati personali pubblicamente accessibili (anagrafica, interessi, relazioni) e
compromette gli account con pochissimi tentativi. È a questo livello che
l'\emph{Open Source Intelligence} (OSINT) diventa decisiva.

Questo lavoro analizza il legame tra informazioni personali esposte online e
sicurezza effettiva delle credenziali: da un lato quantifica, tramite metriche
entropiche, la perdita di resistenza quando i dati personali entrano
nell'attacco; dall'altro descrive le pipeline di raccolta OSINT e generazione
di dizionari mirati, per ricavarne contromisure concrete lato server e
lato utente.
